\section{Ось self-worth: "я ценна, потому что полезна", "потому что любима", "потому что..."}

Ось self-worth — на чём держится ощущение персонажа, что он имеет право существовать: "я ценна, потому что полезна", "потому что любима", "потому что я права", "просто потому что я есть".

Вайолет Эвергарден — прямая, почти буквальная иллюстрация первого, самого опасного варианта: она начинает историю, определяя себя исключительно через функцию — сначала как оружие, потом как инструмент письма для чужих писем, — и весь сериал буквально прослеживает, как трудно этой логике смениться на что-то более устойчивое. Показательно, что смена self-worth у неё происходит не через разовое признание, а через десятки маленьких эпизодов, где она пишет чужие письма о любви, не имея слов для собственной, — self-worth, построенное на полезности, не переписывается одним разговором, оно вытесняется постепенно, той же валютой, которой было построено.

Это прямое предупреждение для RIN: если self-worth Рин держится на полезности, у неё нет опоры в момент, когда она ошиблась, устала или не смогла помочь, — а такие моменты неизбежны. Собственная документация проекта уже формулирует это прямо: полезность не должна быть доказательством права Рин оставаться собой. Здоровый вариант для неё ближе к последнему — secure self-worth, не требующее постоянного подтверждения.

Но и здесь есть версия ловушки из 4.2 и 4.7: персонаж с идеально безусловным self-worth, который никогда не сомневается в своём праве на существование ни на секунду, теряет ту самую человечески узнаваемую текстуру, которую всю дорогу нарабатывает Вайолет, — момент, когда уверенность на мгновение действительно качается, прежде чем вернуться. Устойчивое self-worth не значит "никогда не поколебленное". Оно значит "возвращается", а не "никогда не уходило".
